\documentclass[12pt,letterpaper]{article}
\usepackage[utf8]{inputenc}
\usepackage[T1]{fontenc}
\usepackage{times}
\usepackage[margin=0.5in,top=0.4in,bottom=0.4in]{geometry}
\usepackage{amsmath,amssymb,mathrsfs,braket}
\DeclareMathOperator{\tr}{tr}
\DeclareMathOperator{\rank}{rank}
\DeclareMathOperator{\diag}{diag}
\usepackage{hyperref}
\usepackage{xcolor}
\usepackage{titlesec}
\hypersetup{colorlinks=true,urlcolor=blue!60!black}
\setlength{\parindent}{0pt}
\setlength{\parskip}{3pt}
\titleformat{name=\section,numberless}{\large\bfseries\color{blue!60!black}}{}{0pt}{}[\vspace{-2pt}\rule{\linewidth}{0.6pt}]
\titlespacing{\section}{0pt}{8pt}{4pt}
\pagestyle{plain}
\begin{document}
{\footnotesize\textbf{Quantum Computing Digest --- 2026-11-13} \hfill 8 papers}

{\scriptsize Full author abstracts; italic lines are generated one-line summaries}
\section*{Machine Learning for Quantum Control and Pulse Optimization}
{\small\textbf{A Hardware-Trained Generative Quantum Eigensolver for Sample-Based Quantum Diagonalization of Organotin Hydroxides up to 70 Qubits}} {\scriptsize[\href{https://arxiv.org/abs/2611.13755}{arxiv}]}\nopagebreak\\
{\scriptsize Kunal Kumar, Kimberlee Keithley, Toshiya Iwamori, Shu Kanno, Ikko Hamamura, Takao Kobayashi, Akira Koseki, Qi Gao, Naoki Yamamoto, Hiroshi Horii, Yuri Kobayashi, Borja Peropadre, Kouhei Nakaji, Toru Imai, Takayuki Osogami}\nopagebreak\\[2pt]
{\footnotesize\textit{Hardware-trained circuit generators improve sample-based diagonalization for organotin models up to 70 qubits.}\par}\nopagebreak
{\footnotesize Quantum hardware can provide feedback for training circuit generators beyond the reach of exact classical simulation. We train generative quantum eigensolver (GQE) policies, transformers conditioned on molecular descriptors, to generate the state-preparation circuits for sample-based quantum diagonalization (SQD). Training runs first on a noiseless simulator and then with the IBM Heron processor ibm\_kobe inside the loop, at 46 and 70 qubits ($1.2 \times 10^8$ shots for the four main training runs). A classical CCSD(T) reference energy of the training molecules enters only an admission gate and two auxiliary reward terms. During hardware training, a reward penalty prevents collapse to circuits that sample only the Hartree-Fock determinant. Policies trained on transition-metal complexes are applied without retraining to Sn(OH)$_3$-Me and its hydrates, minimal models of organotin photoresists for extreme-ultraviolet lithography, from 32 to 70 qubits. After one round of hardware sampling, followed by classical ranking and recovery of the cached counts with the molecular Hamiltonian and no further QPU calls, 27 of 28 runs at 46-70 qubits are within 15 mHa of the density-matrix renormalization group (DMRG) reference (median 4.5 mHa), and in all eight runs with an untrained control the trained policy is more accurate. At 56-70 qubits the selected circuits use 34-72 times fewer two-qubit gates than the local unitary cluster Jastrow (LUCJ) circuit transpiled for the same device. These results demonstrate that circuit generators can be trained with quantum hardware at 70 qubits.\par}
\medskip
{\small\textbf{Training quantum neural networks one job at a time}} {\scriptsize[\href{https://arxiv.org/abs/2611.16365}{arxiv}]}\nopagebreak\\
{\scriptsize Muhammad Faryad, Rabab Alomairy}\nopagebreak\\[2pt]
{\footnotesize\textit{Batched parameter-shift updates learn more per cloud job than simultaneous-perturbation training.}\par}\nopagebreak
{\footnotesize When a quantum processor is reached through a cloud queue, a variational algorithm pays for every round trip, and a single round trip can carry many circuits. The classical optimizer therefore decides how much is learned per job, not only per shot. We study this for a quantum neural network classifier. We compare optimizers that are given the same number of jobs and the same shots per job, and that differ only in how they fill each job. This comparison treats hybrid quantum-classical training as an execution-efficiency problem: how to obtain useful optimization progress from each costly interaction with a quantum processor. The resulting perspective connects optimizer design to overhead amortization and efficient use of accelerator resources. A direct port of a batched pattern search that worked for the quantum approximate optimization algorithm stalls once the model has a few dozen parameters. What does transfer is the principle behind it: fill each job with a batch of parameter vectors that is fixed before the job is sent. For circuits built from Pauli rotations the natural batch is the parameter-shift stencil. We show that its gradient error falls as the inverse of the shots per job, while the error of simultaneous-perturbation methods (SPSA) has a floor that no number of shots removes. In simulation, batched parameter-shift training needs between a fifth and a half of the jobs SPSA needs, for 24 to 48 parameters, and the ordering holds at two noise levels. A new structure-aware stencil search reads the same kind of rows as sinusoids and also beats SPSA. On IBM's ibm\_fez processor, in three interleaved runs at six and eight qubits, the batched optimizer reached SPSA's final training loss within 4 to 20 of its 50 jobs. Its test accuracy was clearly higher at eight qubits; at six qubits the difference was not significant. All hardware numbers can be re-derived from the raw shots.\par}
\medskip
\section*{Machine Learning for Error Correction, Mitigation and Fault Tolerance}
{\small\textbf{Improved Low-Cost Quantum Error Mitigation with Benchmarked-Noise Zero-Noise Extrapolation (bnZNE)}} {\scriptsize[\href{https://arxiv.org/abs/2611.16971}{arxiv}]}\nopagebreak\\
{\scriptsize Joseph Harris, Kevin Lively, Peter Ken Schuhmacher}\nopagebreak\\[2pt]
{\footnotesize\textit{Benchmarked-noise extrapolation improves zero-noise estimates at similar measurement cost.}\par}\nopagebreak
{\footnotesize We theoretically extend and experimentally demonstrate benchmarked-noise zero-noise extrapolation (bnZNE), a simple and low-cost method for quantum error mitigation which aims to outperform the highly popular ZNE method for an equal total shot cost by improving the classical extrapolation step. We recently proposed bnZNE as a new error mitigation method (Harris et al., arXiv:2603.10224 (2026)) and showed that it can mitigate noise to the same order as standard ZNE under gate-wise depolarizing noise. Here, we introduce randomized bnZNE which addresses the case of running multiple random circuits per noise level, allowing for the incorporation of noise-tailoring techniques like Pauli twirling and readout error mitigation, and test the method on the IBM and IQM hardware for a plethora of large-scale applications, observing bnZNE to consistently produce more accurate results than standard ZNE. Our approach enhances the extrapolation process to the zero-noise limit by using classically-simulable and structure-mirroring benchmark circuits, generated from the standard application circuit, to sensitively measure the effective error rate at each scaled noise level beyond simple integer values. We give a new method for generating tailored benchmark circuits for the IQM superconducting devices and test them experimentally. Our results demonstrate that bnZNE presents a versatile alternative to ZNE which can enable more accurate results at utility scales with an equivalent cost.\par}
\medskip
{\small\textbf{Scaling P\textbackslash{}\textasciicircum{}elagos: Multi-GPU Orchestration for Relay Belief-Propagation Decoding}} {\scriptsize[\href{https://arxiv.org/abs/2611.16285}{arxiv}]}\nopagebreak\\
{\scriptsize Xulong Tang, Dionysios Diamantopoulos, Tristan M\"uller, Apo Kayi, Nelson Mimura Gonzalez, Muhammad Osama, Seetharami Seelam, Haralampos Pozidis, Blake R. Johnson, Yasuko Eckert}\nopagebreak\\[2pt]
{\footnotesize\textit{Multi-GPU orchestration speeds belief-propagation decoding for large error-correction simulations.}\par}\nopagebreak
{\footnotesize Fault-tolerant quantum computing requires quantum error-correction (QEC) decoders that achieve low logical error rates under real-time operating constraints. These constraints motivate heuristic decoding algorithms whose error-correction performance must be evaluated through large-scale simulation. High-throughput decoding accelerates this evaluation and the exploration of decoder parameters. Using multiple graphics processing units (GPUs) is a natural way to scale these simulations, but naively distributing syndrome batches introduces host-side preprocessing and output-collection overheads that limit scaling. To address these limitations, we extend the P\textbackslash{}\textasciicircum{}elagos relay belief-propagation (Relay-BP) decoder with process-isolated orchestration, parallel host preprocessing, and shared-memory output. We evaluate both one-gross ($[[144,12,12]]$) and two-gross ($[[288,12,18]]$) codes in a multi-GPU environment. Specifically, for the one-gross circuit at $R=50$, eight AMD Instinct MI300X GPUs and eight AMD Instinct MI355X GPUs achieve optimized end-to-end speedups of $5.68\times$ and $6.23\times$, respectively. For the two-gross circuit, the corresponding speedups are $6.92\times$ and $7.01\times$. When throughput is calculated using GPU decoding time, the corresponding two-gross speedups are $7.58\times$ and $7.45\times$, indicating near-linear GPU scaling. We deploy P\textbackslash{}\textasciicircum{}elagos on Lux, Oak Ridge National Laboratory's AMD Instinct MI355X-based supercomputer, as a pathfinding application for high-throughput QEC simulation. This deployment provides a concrete quantum-computing use case for the AI + HPC infrastructure supporting the Genesis Mission's scientific-discovery goals. These end-to-end measurements demonstrate a practical path to accelerating large-scale QEC simulation through efficient multi-GPU decoding.\par}
\medskip
\section*{Notable Quantum Computing Research}
{\small\textbf{An Architecture for Real-Time Quantum Error Correction with Neutral Atoms}} {\scriptsize[\href{https://arxiv.org/abs/2611.15104}{arxiv}]}\nopagebreak\\
{\scriptsize X. Wang, Y. Jing, Z. Yu, S. Wang, D. Niu, Y. Liu, T. Xiao, A. Yang, J. Han, Y. Wu, X. Lv}\nopagebreak\\[2pt]
{\footnotesize\textit{A heterogeneous pipeline overlaps readout and decoding to cut neutral-atom error-correction latency.}\par}\nopagebreak
{\footnotesize The latency of quantum measurement and classical decoding constrains the execution of adaptive fault-tolerant circuits. In neutral-atom processors, this requirement couples spatially resolved fluorescence detection to classical inference. Here we implement cascaded readout and decoding in a heterogeneous architecture that overlaps acquisition and processing at two spatial scales. Field-programmable gate arrays process camera rows as they arrive, while graphics processors decode an acquired logical block during acquisition of the next. This heterogeneous architecture offers flexibility in the choice of error-correcting codes and decoding algorithms. Direct communication preserves microsecond-scale transfer times for kilobyte payloads and supports interchangeable decoding algorithms. Hardware characterization and decoding benchmarks using circuit-level simulated syndromes yield a submillisecond post-exposure processing latency for selected small-code workloads ($d = 3, 5, 7$). By enabling logical blocks to be decoded before full-frame readout is complete, the architecture implements low-latency syndrome decoding, realizing real-time quantum information processing as required for fault-tolerant computation in atom arrays.\par}
\medskip
{\small\textbf{Scalable Compilation of Resource-Efficient Tapers for High-Confidence Quantum Phase Estimation}} {\scriptsize[\href{https://arxiv.org/abs/2611.15076}{arxiv}]}\nopagebreak\\
{\scriptsize David Raveh, Andrew Sornborger, Yi\u{g}it Suba\c{s}\textbackslash{}i, Lukasz Cincio}\nopagebreak\\[2pt]
{\footnotesize\textit{Tensor-network compilation builds shallow taper-state circuits for high-confidence quantum phase estimation.}\par}\nopagebreak
{\footnotesize Quantum phase estimation is a key component of many quantum algorithms. Recently, high-confidence, tapered QPE (tQPE) has been introduced, significantly reducing the number of ancilla qubits necessary for a given failure probability, $\epsilon$. We consider tQPE with $p = l+m$ ancillas, where $l$ represents the number of ancillas necessary to estimate a phase to a given resolution and $m$ represents the number of ''confidence'' qubits added to ensure a desired success probability. With tQPE the failure probability decays double exponentially in $m$, implying that only a few confidence qubits suffice for virtually any problem of interest. A key challenge for tQPE is the construction of the quantum circuit for preparing the taper state on the ancilla register. Here, we develop scalable tensor network methods (previous work was limited to ($l + m = 24$)) giving highly optimized state preparation circuits for optimally concentrated DPSS (Slepian) tapers. We provide an accompanying database of resource-efficient state preparation circuits for \$l = 1, \textbackslash{}dots, 50\$ and \$m = 1, \textbackslash{}dots, 4\$. Our circuits are shallow enough for a demonstration on near-term devices for small values of $l$. We run tQPE circuits for a 1-qubit system initialized in the ground state on Quantinuum's H2-2 quantum computer, demonstrating a $3\times$ to $18\times$ reduction in failure probability for $l=1,\ldots,10$ on a noisy device.\par}
\medskip
{\small\textbf{Reducing TLS Losses in Superconducting Circuits Using Solid Neon}} {\scriptsize[\href{https://arxiv.org/abs/2611.15021}{arxiv}]}\nopagebreak\\
{\scriptsize A. Palmer, K. Matkovic, P. Russell, E. Helgemo, L. Delventhal, Y. Ashlea Alava, K. Zuo, S. N. Coppersmith, M. C. Cassidy}\nopagebreak\\[2pt]
{\footnotesize\textit{Solid neon treatment reduces losses in resonators, likely by lowering two-level-system dissipation.}\par}\nopagebreak
{\footnotesize Two level systems (TLSs) are a major source of loss in superconducting devices, limiting the performance of qubits, detectors, and resonators. These defects arise from a range of sources, including surface oxides, crystal defects, as well as fabrication residues which are often difficult to remove. Using NbTiN coplanar waveguide resonators, we show that depositing and subsequently removing a solid neon film can significantly improve device performance. Analysis of the temperature and power dependence indicates that this enhancement arises from reduced TLS dissipation, which we attribute to the redistribution and removal of surface contaminants during neon deposition and evaporation. Atomic force microscopy images of contaminated hexagonal boron nitride flakes before and after neon exposure show that neon treatment can redistribute and remove surface residues. These results demonstrate a simple, non-destructive approach for mitigating TLS losses in superconducting devices and other sensitive materials that cannot withstand aggressive cleaning processes.\par}
\medskip
{\small\textbf{Persistent-mediator architecture for universal fault-tolerant quantum computation with color codes}} {\scriptsize[\href{https://arxiv.org/abs/2611.16691}{arxiv}]}\nopagebreak\\
{\scriptsize Zhi-Cheng He, Jingbo Wang, Dong E. Liu}\nopagebreak\\[2pt]
{\footnotesize\textit{A persistent 3D mediator enables repeated logical T gates while keeping 2D data encoded.}\par}\nopagebreak
{\footnotesize Two-dimensional color codes provide transversal Clifford gates, whereas suitable three-dimensional codes provide a transversal $T$ gate. Representative cross-dimensional schemes either switch the computational encoding or consume a 3D encoded state. Reusing the same protected 3D block would remove repeated auxiliary-state preparation and terminal readout, raising a central question: can the 3D module remain encoded and protected while mediating successive $T$ gates on a separate 2D register? Here we introduce a cross-dimensional persistent-mediator architecture that retains the 2D data register and 3D mediator as distinct encoded modules at each cycle boundary, coupled only through a transient protected subsystem. We construct a membrane-like joint-logical measurement and combine it with a matched cross-dimensional $CZ$, complementary string-like surgery, and feed-forward to close a deterministic logical $T$-gate cycle. The cycle avoids post-selection and destructive mediator readout and restores the 3D block to its initial logical state. Stabilizer analysis establishes the logical action and preserved spatial distance. Phenomenological and circuit-level simulations of a Clifford surrogate show finite-distance scaling close to the matched 3D memory benchmark, while sequential tests show approximately linear logical-error accumulation in the low-error regime. Persistent mediation enables repeated access to higher-dimensional non-Clifford capability while keeping the computational register in its 2D encoding.\par}
\medskip
\end{document}